\section{符号说明}

下表收录贯穿全文、直接进入核心方程与判据的符号。

\begin{longtable}{clcl}
\caption*{主要符号说明}\\
\toprule
符号 & 含义 & 单位 & 取值或范围 \\
\midrule
\endfirsthead
\toprule
符号 & 含义 & 单位 & 取值或范围 \\
\midrule
\endhead
\bottomrule
\endfoot
$r$ & 到药材轴心的径向距离 & m & $[0,R]$ \\
$t$ & 时间 & s & $[0,t_{\rm end}]$ \\
$T(r,t)$ & 药材内部温度 & $^\circ$C & $[28,\,50.2]$ \\
$C(r,t)$ & 药材内部干基含水率 & kg/kg & $[0.0526,\,2.55]$ \\
$T_K$ & 绝对温度 $T+273.15$ & K & $[301.15,\,323.32]$ \\
$w$ & 湿基水分分率 $C/(C+1)$ & --- & $[0.050,\,0.718]$ \\
$\rho$ & 药材密度 & kg$\cdot$m$^{-3}$ & 随物性关系而定 \\
$c_p$ & 定压比热容 & J$\cdot$kg$^{-1}\cdot$K$^{-1}$ & 随物性关系而定 \\
$k$ & 导热系数 & W$\cdot$m$^{-1}\cdot$K$^{-1}$ & 随物性关系而定 \\
$D$ & 水分有效扩散系数 & m$^2\cdot$s$^{-1}$ & $1.6\times10^{-10}$--$1.4\times10^{-8}$ \\
$h$ & 表面对流换热系数 & W$\cdot$m$^{-2}\cdot$K$^{-1}$ & $25.0$ \\
$h_m$ & 表面对流传质系数 & m$\cdot$s$^{-1}$ & $8\times10^{-7}$ \\
$T_{\rm air}(t)$ & 烘房空气温度 & $^\circ$C & $[28.000,\,50.246]$ \\
$C_{\rm air}(t)$ & 烘房空气含湿量 & kg/kg & $[0.01963,\,0.05025]$ \\
$R(t)$ & 药材半径 & m & $0.0200\to0.01198$ \\
$\dot R(t)$ & 半径收缩速率 & m$\cdot$s$^{-1}$ & $\le 0$ \\
$\xi$ & 归一化径向坐标 $r/R(t)$ & --- & $[0,1]$ \\
$\eta$ & 归一化物质坐标 $(r/R(t))^2$ & --- & $[0,1]$ \\
$v_s$ & 固相仿射速度 $r\dot R/R$ & m$\cdot$s$^{-1}$ & $\le 0$ \\
$\chi$ & 参考系开关（$1$ 物质坐标主模型，$0$ 保留表观对流） & --- & $\{0,1\}$ \\
$C_{\rm th}$ & 达标含水率阈值 & kg/kg & $0.15$ \\
$t_{\rm end}$ & 烘干结束时间 & h & 见各问结果 \\
$\alpha$ & 热扩散率 $k/(\rho c_p)$ & m$^2\cdot$s$^{-1}$ & $1.689\times10^{-7}$ \\
$N$ & 径向网格区间数 & --- & $160,\,1280,\,2560$ \\
$\Delta t$ & 时间步长 & s & $0.125,\,0.25$ \\
\end{longtable}
\addtocounter{table}{-1}
